% GENERATED FILE. Edit canonical lessons, then run npm run generate:books.
% canonical-source-hash: fnv1a64:7b44d778c1697dd8
% canonical-lessons: BN-C332-sutorang, BN-C332-tachhara, BN-C332-othoch, BN-C332-oboshyo, BN-C332-khata

\chapter{One Thing Leads to Another}
\label{ch:bn-one-thing-leads-to-another}
\hlchaptermodality{332}

\begin{chapteropening}
\textbf{By the end of this chapter:} \emph{I can link one reason to the next and say what follows.}

Complete the last lesson of chapter 332: I can link one reason to the next and say what follows.
\end{chapteropening}

\section[\texorpdfstring{\bn{সুতরাং} (sutôrāng)}{সুতরাং (sutôrāng)}]{\bn{সুতরাং} (sutôrāng) — therefore; and so}

\label{lesson:BN-C332-sutorang}



\begin{quote}
Before the new one: say the Bengali for a financial loss, then the Bengali for useful.
\end{quote}

\subsection*{You'll want to know: \bn{সুতরাং}}
\textbf{\bn{সুতরাং}} — \emph{sutôrāng} — \textquotedblleft{}therefore; and so\textquotedblright{}.

\begin{quote}
\bn{বাস} \bn{নেই}, \bn{সুতরাং} \bn{আমরা} \bn{হেঁটে} \bn{যাব।}
\end{quote}

\begin{quote}
\emph{bās nei, sutôrāng āmrā hẽṭe jābô.}
\end{quote}

\begin{quote}
There is no bus, so we will walk.
\end{quote}

\begin{cousinweb}
\textbf{\bn{সুতরাং}} (\emph{sutôrāng}) is Sanskrit \emph{sutarām}, 'all the more' (Monier-Williams), with its final m heard as ng. The sense 'therefore' is not one the dictionary gives; in Bengali the word draws the conclusion from what was said before.
\end{cousinweb}

\begin{grammarlens}[title={What you've built}]
One more word for giving reasons and opinions.
\end{grammarlens}

\subsection*{Your turn}
\begin{itemize}
\raggedright
  \item \emph{Say it:} \emph{sutôrāng}
  \item \emph{Say it:} \emph{sutôrāng}, once more
  \item \emph{Say it:} \emph{bās nei, sutôrāng āmrā hẽṭe jābô.} — There is no bus, so we will walk.
  \item \emph{Recall:} say the Bengali for a financial loss, then the Bengali for useful, then say \emph{sutôrāng} again
\end{itemize}

\begin{tcolorbox}[breakable,colback=teal!4,colframe=teal!35!black,title={Before you move on}]
What does \bn{সুতরাং} mean? (\textquotedblleft{}Therefore; and so\textquotedblright{}.) Say it once more.
\end{tcolorbox}
\section[\texorpdfstring{\bn{তাছাড়া} (tāchhāṛā)}{তাছাড়া (tāchhāṛā)}]{\bn{তাছাড়া} (tāchhāṛā) — besides; apart from that}

\label{lesson:BN-C332-tachhara}



\begin{quote}
Before the new one: say the Bengali for useful, then the Bengali for therefore.
\end{quote}

\subsection*{You'll want to know: \bn{তাছাড়া}}
\textbf{\bn{তাছাড়া}} — \emph{tāchhāṛā} — \textquotedblleft{}besides; apart from that\textquotedblright{}.

\begin{quote}
\bn{আজ} \bn{আমি} \bn{ক্লান্ত}, \bn{তাছাড়া} \bn{আমার} \bn{অনেক} \bn{কাজ} \bn{আছে।}
\end{quote}

\begin{quote}
\emph{āj āmi klāntô, tāchhāṛā āmār ônek kāj āchhe.}
\end{quote}

\begin{quote}
I'm tired today, and besides, I have a lot of work.
\end{quote}

\begin{cousinweb}
\textbf{\bn{তাছাড়া}} (\emph{tāchhāṛā}) is \bn{তা} 'that' with \bn{ছাড়া} 'apart from' (CDIAL 5612 and 4998 list Bengali tā and chāṛā): 'apart from that'. It adds a second reason to the first.
\end{cousinweb}

\begin{grammarlens}[title={What you've built}]
One more word for giving reasons and opinions.
\end{grammarlens}

\subsection*{Your turn}
\begin{itemize}
\raggedright
  \item \emph{Say it:} \emph{tāchhāṛā}
  \item \emph{Say it:} \emph{tāchhāṛā}, once more
  \item \emph{Say it:} \emph{āj āmi klāntô, tāchhāṛā āmār ônek kāj āchhe.} — I'm tired today, and besides, I have a lot of work.
  \item \emph{Recall:} say the Bengali for useful, then the Bengali for therefore, then say \emph{tāchhāṛā} again
\end{itemize}

\begin{tcolorbox}[breakable,colback=teal!4,colframe=teal!35!black,title={Before you move on}]
What does \bn{তাছাড়া} mean? (\textquotedblleft{}Besides; apart from that\textquotedblright{}.) Say it once more.
\end{tcolorbox}
\section[\texorpdfstring{\bn{অথচ} (ôthôch)}{অথচ (ôthôch)}]{\bn{অথচ} (ôthôch) — and yet; even so}

\label{lesson:BN-C332-othoch}



\begin{quote}
Before the new one: say the Bengali for therefore, then the Bengali for besides.
\end{quote}

\subsection*{You'll want to know: \bn{অথচ}}
\textbf{\bn{অথচ}} — \emph{ôthôch} — \textquotedblleft{}and yet; even so\textquotedblright{}.

\begin{quote}
\bn{মীরা} \bn{অনেক} \bn{পড়ল}, \bn{অথচ} \bn{পাস} \bn{করল} \bn{না।}
\end{quote}

\begin{quote}
\emph{mīrā ônek pôṛlô, ôthôch pās kôrlô nā.}
\end{quote}

\begin{quote}
Mira studied a lot, and yet she did not pass.
\end{quote}

\begin{cousinweb}
\textbf{\bn{অথচ}} (\emph{ôthôch}) is Sanskrit \emph{atha ca}, 'and moreover' (Monier-Williams). The Bengali sense, 'and yet', which links a result you would not expect, is not one the dictionary gives.
\end{cousinweb}

\begin{grammarlens}[title={What you've built}]
One more word for giving reasons and opinions.
\end{grammarlens}

\subsection*{Your turn}
\begin{itemize}
\raggedright
  \item \emph{Say it:} \emph{ôthôch}
  \item \emph{Say it:} \emph{ôthôch}, once more
  \item \emph{Say it:} \emph{mīrā ônek pôṛlô, ôthôch pās kôrlô nā.} — Mira studied a lot, and yet she did not pass.
  \item \emph{Recall:} say the Bengali for therefore, then the Bengali for besides, then say \emph{ôthôch} again
\end{itemize}

\begin{tcolorbox}[breakable,colback=teal!4,colframe=teal!35!black,title={Before you move on}]
What does \bn{অথচ} mean? (\textquotedblleft{}And yet; even so\textquotedblright{}.) Say it once more.
\end{tcolorbox}
\section[\texorpdfstring{\bn{অবশ্য} (ôbôshyô)}{অবশ্য (ôbôshyô)}]{\bn{অবশ্য} (ôbôshyô) — of course; mind you}

\label{lesson:BN-C332-oboshyo}



\begin{quote}
Before the new one: say the Bengali for besides, then the Bengali for and yet.
\end{quote}

\subsection*{You'll want to know: \bn{অবশ্য}}
\textbf{\bn{অবশ্য}} — \emph{ôbôshyô} — \textquotedblleft{}of course; mind you\textquotedblright{}.

\begin{quote}
\bn{হোটেলটা} \bn{দামি}, \bn{অবশ্য} \bn{খাবার} \bn{দুর্দান্ত।}
\end{quote}

\begin{quote}
\emph{hoṭelṭā dāmi, ôbôshyô khābār durdāntô.}
\end{quote}

\begin{quote}
The hotel is expensive, though mind you, the food is excellent.
\end{quote}

\begin{cousinweb}
\textbf{\bn{অবশ্য}} (\emph{ôbôshyô}) is Sanskrit \emph{avaśyam}, 'necessarily' (Monier-Williams). Bengali uses it to grant a point while keeping your own: 'of course', 'mind you'.
\end{cousinweb}

\begin{grammarlens}[title={What you've built}]
One more word for giving reasons and opinions.
\end{grammarlens}

\subsection*{Your turn}
\begin{itemize}
\raggedright
  \item \emph{Say it:} \emph{ôbôshyô}
  \item \emph{Say it:} \emph{ôbôshyô}, once more
  \item \emph{Say it:} \emph{hoṭelṭā dāmi, ôbôshyô khābār durdāntô.} — The hotel is expensive, though mind you, the food is excellent.
  \item \emph{Recall:} say the Bengali for besides, then the Bengali for and yet, then say \emph{ôbôshyô} again
\end{itemize}

\begin{tcolorbox}[breakable,colback=teal!4,colframe=teal!35!black,title={Before you move on}]
What does \bn{অবশ্য} mean? (\textquotedblleft{}Of course; mind you\textquotedblright{}.) Say it once more.
\end{tcolorbox}
\section[\texorpdfstring{\bn{খাটা} (khāṭā)}{খাটা (khāṭā)}]{\bn{খাটা} (khāṭā) — to hold good; to apply (of a rule); to toil}

\label{lesson:BN-C332-khata}



\begin{quote}
Before the new one: say the Bengali for and yet, then the Bengali for mind you.
\end{quote}

\subsection*{You'll want to know: \bn{খাটা}}
\textbf{\bn{খাটা}} — \emph{khāṭā} — \textquotedblleft{}to hold good; to apply (of a rule); to toil\textquotedblright{}.

\begin{quote}
\bn{এই} \bn{নিয়মটা} \bn{এখানে} \bn{খাটে} \bn{না।}
\end{quote}

\begin{quote}
\emph{ei niyômṭā ekhāne khāṭe nā.}
\end{quote}

\begin{quote}
This rule doesn't apply here.
\end{quote}

Past (he/she …): \bn{সেই} \bn{নিয়মটা} \bn{খাটল} \bn{না}, \emph{sei niyômṭā khāṭlô nā}.

\begin{cousinweb}
\textbf{\bn{খাটা}} (\emph{khāṭā}) is listed by CDIAL 3779 (\emph{khaṭṭáyati}) with the senses 'to work' and 'to apply'. A person \bn{খাটে} when working hard; a rule or a reason \bn{খাটে} when it holds in a case.
\end{cousinweb}

\begin{grammarlens}[title={What you've built}]
One more word for giving reasons and opinions.
\end{grammarlens}

\subsection*{Your turn}
\begin{itemize}
\raggedright
  \item \emph{Say it:} \emph{khāṭā}
  \item \emph{Say it:} \emph{khāṭā}, once more
  \item \emph{Say it:} \emph{ei niyômṭā ekhāne khāṭe nā.} — This rule doesn't apply here.
  \item \emph{Recall:} say the Bengali for and yet, then the Bengali for mind you, then say \emph{khāṭā} again
\end{itemize}

\begin{tcolorbox}[breakable,colback=teal!4,colframe=teal!35!black,title={Before you move on}]
What does \bn{খাটা} mean? (\textquotedblleft{}To hold good; to apply (of a rule); to toil\textquotedblright{}.) Say it once more.
\end{tcolorbox}
